% =====================================================================
%  Chapitre 6 - Sprint 3 : demande, devis, mission et suivi en direct
% =====================================================================
\chapter{Sprint 3 : demande, devis, mission et suivi en direct}
\label{chap:sprint3}

\section*{Introduction}
\addcontentsline{toc}{section}{Introduction}

Les deux sprints precedents ont pose les socles : qui utilise la
plateforme, et avec quelles ressources. Celui-ci aborde \essentiel{le
coeur de l'activite} : la chaine qui va de la demande du client a la
livraison de sa marchandise, en passant par le devis, l'affectation et le
suivi. C'est le sprint le plus dense du projet, et il mobilise deux
services : celui des missions et le moteur d'estimation tarifaire.

% =====================================================================
\section{Planification du sprint}

\subsection{Objectifs}

\begin{enumerate}
  \item Ouvrir au client un canal autonome de demande de transport.
  \item Estimer le prix d'un trajet avant meme que le devis ne soit
    etabli.
  \item Formaliser le cycle du devis, de l'envoi a l'acceptation ou au
    refus.
  \item Creer la mission, l'affecter et en suivre l'execution.
  \item Restituer la position du vehicule en direct sur une carte.
\end{enumerate}

\subsection{Carnet du sprint}

\begin{table}[H]
\centering
\caption{Carnet du sprint 3}
\label{tab:backlog-s3}
\small
\begin{tabularx}{\textwidth}{@{}cp{2.3cm}Xcc@{}}
\toprule
\textbf{N\textsuperscript{o}} & \textbf{Acteur} & \textbf{Besoin exprime} &
\textbf{Prio.} & \textbf{Pts} \\ \midrule
1 & Client & Deposer une demande de transport & Must & 8 \\
2 & Client & Connaitre une fourchette de prix avant de valider & Should & 8 \\
3 & Responsable & Etablir et envoyer un devis & Must & 8 \\
4 & Client & Accepter ou refuser un devis recu & Must & 5 \\
5 & Responsable & Affecter un vehicule et un chauffeur & Must & 5 \\
6 & Chauffeur & Transmettre sa position depuis l'application mobile & Should & 8 \\
7 & Client & Suivre sa mission en direct sur une carte & Should & 8 \\
8 & Responsable & Cloturer ou annuler une mission & Must & 3 \\
9 & Client & Ouvrir un ticket de support sur une mission & Could & 3 \\
10 & Responsable & Consulter l'historique d'un trajet & Could & 3 \\
\midrule
\multicolumn{4}{@{}r}{\textbf{Total}} & \textbf{59} \\
\bottomrule
\end{tabularx}
\end{table}

% =====================================================================
\section{Elaboration technique}

\subsection{Modelisation fonctionnelle}

\figReelle{mission-usecase.png}
  {Cas d'utilisation du service des missions}
  {fig:uc-mission}{0.54}

Le diagramme de la figure~\ref{fig:uc-mission} organise le service en
quatre domaines : la demande et le devis, l'execution des missions, le
suivi en direct et l'assistance. Trois cas se detachent a droite de la
frontiere, atteints par inclusion ou par extension : l'estimation du
prix, la creation de la mission et l'affectation des ressources.

\paragraph{Une regle metier structurante.}
\essentiel{Deposer une demande n'est pas commander une mission.} La
demande de transport est un \exemple{devis} : elle vit sa propre vie,
recoit un prix, puis est acceptee ou refusee. \essentiel{La mission n'est
creee qu'a l'acceptation}, ce que traduit la relation d'inclusion entre
\exemple{accepter un devis} et \exemple{creer une mission}.

Cette regle a une consequence concrete sur l'interface de programmation :
le point d'entree qui permettait a un client de creer directement une
mission a ete \essentiel{retire}. Le laisser aurait ouvert un chemin
parallele contournant le devis, et donc la tarification.

\begin{table}[H]
\centering
\caption{Description textuelle du cas « accepter un devis »}
\label{tab:desc-devis}
\small
\begin{tabularx}{\textwidth}{@{}p{2.5cm}X@{}}
\toprule
\textbf{Rubrique} & \textbf{Contenu} \\ \midrule
Acteur & Client \\
Precondition & La demande porte le statut QUOTATION\_SENT et un prix
  final a ete fixe \\
Scenario nominal & Le client consulte le devis, l'accepte ; le systeme
  cree la mission au statut EN\_ATTENTE, rattache la demande a la mission,
  fait passer la demande au statut ACCEPTED et publie un evenement sur
  le bus \\
Alternatives & Refus : la demande passe au statut REJECTED avec le motif
  saisi, aucune mission n'est creee. Devis deja traite : l'action est
  refusee \\
Postcondition & La mission existe et attend l'affectation d'un vehicule
  et d'un chauffeur \\
\bottomrule
\end{tabularx}
\end{table}

\subsection{Conception statique}

\figReelle{mission-class-domain.png}
  {Modele du domaine du service des missions}
  {fig:dom-mission}{0.30}

Le modele distingue nettement \exemple{DemandeTransport}, qui porte la
negociation (villes, marchandise, prix estime, fourchette, prix final,
methode de tarification), et \exemple{Mission}, qui porte l'execution
(vehicule, chauffeur, kilometrages, cout reel). Le lien entre les deux
est \essentiel{optionnel et oriente} : une demande peut exister sans
mission, jamais l'inverse.

\paragraph{Trois methodes de tarification.}
Le prix final s'accompagne d'une methode indiquant comment il a ete
obtenu : propose par le moteur et retenu tel quel, modifie, ou fixe
manuellement. \essentiel{Cette information est conservee}, ce qui permet
de mesurer a posteriori dans quelle mesure l'entreprise suit les
recommandations du modele.

\subsection{Conception dynamique}

\figReelle{mission-seq-demande-transport.png}
  {Sequence : depot d'une demande de transport}
  {fig:seq-demande}{0.54}

\figReelle{mission-seq-devis.png}
  {Sequence : etablissement puis acceptation ou refus d'un devis}
  {fig:seq-devis}{0.54}

La premiere sequence montre l'appel au moteur d'estimation \essentiel{avant}
l'enregistrement : le client voit la fourchette tarifaire au moment ou il
saisit sa demande. La seconde, figure~\ref{fig:seq-devis}, comporte un
bloc alternatif : l'acceptation cree la mission et publie un evenement,
le refus se contente d'enregistrer le motif.

\figReelle{mission-seq-suivi-direct.png}
  {Sequence : suivi d'une mission en direct}
  {fig:seq-suivi}{0.54}

Le suivi, figure~\ref{fig:seq-suivi}, repose sur une boucle : tant que la
mission est en cours, l'application du chauffeur transmet sa position, que
le client consulte. \essentiel{Ces positions ne transitent pas par la
passerelle} : l'application mobile s'adresse directement au service des
missions. Le detail figure a la section~\ref{sec:archi-physique}.

\subsection{Integration avec les services d'aide a la decision}

Le moteur d'estimation est un service Python autonome. Il charge en
memoire un modele de gradient boosting entraine sur un historique de
trajets, ainsi qu'une matrice de distances entre les villes desservies.
Sur reception d'une demande, il calcule la distance, applique le modele et
renvoie un prix encadre d'une fourchette basse et haute.

\figReelle{mission-communication.png}
  {Echanges du service des missions}
  {fig:com-mission}{0.60}

La figure~\ref{fig:com-mission} montre les deux natures d'echange du
service. L'appel au moteur d'estimation est \essentiel{synchrone} : le
client attend le prix pour decider. La demande d'optimisation d'une
tournee est \essentiel{asynchrone}, publiee sur le bus et traitee en
differe. Le critere est simple : \essentiel{on attend ce dont on a besoin
sur le champ, on delegue le reste}.

% =====================================================================
\section{Realisation et tests de validation}

\figDeux{s3-portail-reservation-1.jpg}{s3-estimation-tarifaire-1.jpg}
  {Portail client : formulaire de demande et fourchette tarifaire
   proposee}
  {fig:ecrans-demande}{0.46}

\figTrois{s3-demandes-transport.jpg}{s3-suivi-devis-client.jpg}{s3-carte-suivi.jpg}
  {Traitement d'une demande cote back-office, suivi du devis cote client,
   et carte de suivi d'une mission en cours}
  {fig:ecrans-mission}{0.40}

\begin{table}[H]
\centering
\caption{Tests de validation du sprint 3}
\label{tab:tests-s3}
\small
\begin{tabularx}{\textwidth}{@{}cXp{2.9cm}c@{}}
\toprule
\textbf{N\textsuperscript{o}} & \textbf{Cas teste} & \textbf{Attendu} &
\textbf{Etat} \\ \midrule
T1 & Depot d'une demande & Statut PENDING\_QUOTATION, prix estime affiche & Reussi \\
T2 & Envoi d'un devis modifie & Methode de tarification MODIFIED conservee & Reussi \\
T3 & Acceptation d'un devis & Mission creee, demande rattachee & Reussi \\
T4 & Refus d'un devis & Motif enregistre, aucune mission creee & Reussi \\
T5 & Position transmise par le chauffeur & Point visible sur la carte du client & Reussi \\
T6 & Trajet inconnu de la matrice & Estimation refusee, saisie manuelle possible & Reussi \\
\bottomrule
\end{tabularx}
\end{table}

% =====================================================================
\section*{Conclusion}
\addcontentsline{toc}{section}{Conclusion}

Ce sprint livre la chaine commerciale complete, de la demande au suivi en
direct. Les cinquante-neuf points du carnet ont ete realises, ce qui en
fait le sprint le plus lourd du projet.

Deux decisions en ressortent. La premiere est metier : \essentiel{la
reservation est un devis, la mission n'existe qu'apres acceptation} ;
elle a conduit a retirer un point d'entree qui permettait de la
contourner. La seconde est une regle d'integration : \essentiel{synchrone
pour ce dont l'utilisateur a besoin sur le champ, asynchrone pour le
reste}, reprise telle quelle par les sprints suivants.

Le sprint suivant traite du maintien en condition operationnelle de la
flotte : maintenance curative, puis anticipation des pannes.
